% !TeX root = ../papers/nc-p3.tex
\section{四次元負成分の楕円点グラフと枠を保つ特殊化}
\label{sec:rectangle-component-point-geometry}
\begin{theorem}\label{thm:negative-component-elliptic-graph}
定理\ref{thm:rectangle-unrestricted-components}の次元4の負成分では、
一般の各辺の点対応は滑らかな楕円四次曲線と四つの被約孤立点の同型グラフである。
この成分のどの窓でも各辺の点対応は正次元である。
\end{theorem}
\begin{proof}
$A_0=\left(\begin{smallmatrix}a&b\\c&d\end{smallmatrix}\right)$,
$A_1=-A_0$, $\lambda=-1$, $D=ad-bc$ とする。第一辺の六関係の左評価は
\[
\begin{pmatrix}
 ax_0&cx_1&x_2&-zx_3\\ bx_0&dx_1&-x_2&x_3\\
 -bx_1&-ax_0&x_3&-x_2\\ -dx_1&-cx_0&-zx_3&x_2\\
 x_2&zx_3&x_0&x_1\\ x_3&x_2&x_1&x_0
\end{pmatrix}.
\]
右評価では各双一次式を転置する。各射影に対して二本の二次式を
$Hx_0^2=Lx_2^2+Mx_3^2$, $Hx_1^2=Nx_2^2+Ox_3^2$ と書く。左では
\[
\begin{split}
 H&=(a-d)D,\quad L=-(b+d)(c+d),\quad M=(c+dz)(bz+d),\\
 N&=-(a+b)(a+c),\qquad O=(a+bz)(az+c).
\end{split}
\]
右では
\[
\begin{split}
 H&=(bz-c)D,\quad L=-(c+d)(az+c),\quad M=(a+c)(c+dz),\\
 N&=-(a+b)(bz+d),\qquad O=(a+bz)(b+d).
\end{split}
\]
いずれも一般点で $HLMNO(LO-MN)\ne0$ なので滑らかな四次完全交叉である。
左では $(x_2=x_3=0,\ x_0^2=x_1^2)$ と
$(x_0=x_1=0,\ x_2^2=zx_3^2)$ の各二点を加える。右では二つの係数 $1,z$ を入れ替える。
この和のイデアルで十五個の四次小行列式を割ると、関数体上で余りは全て零である。
$(a,b,c,d,z)=(5,2,3,7,3)$ では左右とも四アフィンチャートで逆包含も成り立ち、
三次小行列式の共通零点はない。\path{computations/verify_negative_component_point_graph.py} に
六関係の元の切断との照合、関数体での除算、有理特殊点の両方向の検算を与えた。

既知の和は標本の近くで平坦である。包含を与える構造層の全射の核は、平坦性により標本上で消え、
台の射影が閉なので近傍でも消える。三次小行列式の無共通零点性も開条件である。
従って両射影は局所的に階数3の核を取るグラフであって、上の和にスキームとして同型である。
他の辺には捻りの前の両対称性を使う。文字を平行移動する際の
双指標の変化は各文字の対角倍率となる。捻りの後の逆辺は転置に加えて
$\mu(i,j)\mu(j,i)$ が掛かるが、$q(i)=\mu(i,i)$ とすればこれは $q(i)q(j)q(i+j)$ である。
関係は色 $i+j$ ごとに斉次なので逆辺も各文字の基底変更で元の辺の転置に戻る。
各辺の射影的点対応のファイバー次元の上半連続性と成分の稠密表示から最後の主張を得る。
\end{proof}

\begin{lemma}\label{lem:stable-twenty-frames}
辺の六式の零スキームが有限なら長さ20の完全交叉で、二射影は閉埋込みである。
一射影の相異なる点で一つの $r$ 次元射影部分空間に入るものは、
$r=0,1,2$ に応じて高々 $1,4,10$ 個である。
離散付値環上の二つのそのような配置で、特殊点の相異なる点の個数を $n_1,n_2$ とする。
$n_1+n_2\ge32$ なら、一般ファイバーの射影同型は特殊点に可逆に延びる。
この主張の一般ファイバーの同型は、スキーム同型を与える射影線型変換をいう。
\end{lemma}
\begin{proof}
有限なら各部分評価は零点上で階数3でなければならない。階数2以下なら核が射影直線を含む。
可逆三次小行列の開集合で核を解くと射影は閉埋込みである。
六因子の交点数は $(H_1+H_2)^6=20H_1^3H_2^3$ である。
$\mathbb P^r\times\mathbb P^3$ に制限して $r+3$ 本の一般線型結合を取る。
共通基点は有限で、各段階で正次元成分を含む切断を避けられるから、その長さは
$\binom{r+3}{r}$ となり点数の上界を得る。

相対完全交叉は有限平坦である。必要なら付値環を有限拡大して一般点の台を分解する。
重複度込みの長さは各特殊点に特殊化する一般点の重複度の和である
（完備化後の有限自由代数を特殊点の冪等元で分解する）。
任意の $r$ 次元空間に載る特殊点の総重複度は上の上界に $20-n_i$ を足した値以下である。
一般ファイバーの射影変換をSmith標準形で最小指数0にすると、
特殊点で階数 $e<4$ なら核の外に特殊化する一般点の像は全て特殊行列の像に特殊化する。
従って総数20は、次元 $3-e$ の核の重複度と次元 $e-1$ の像の重複度の和以下となる。
ところが異なる点の上界の和は $e=1,3$ で11、$e=2$ で8であり、
二つの重複補正の和は8以下である。矛盾により射影変換は可逆に延びる。一意性は分離性による。
\end{proof}

\begin{proposition}\label{prop:stable-component-markings}
四位置の有限配置がそれぞれ十六個以上の相異なる点を持つ軌跡では、
一般点での位置移動を伴う窓同型、四元群の正則文字枠、および一般点での両対称表示は特殊化する。
特に次元14の正成分のこの部分における全ての窓は両対称表示を持ち、
定理\ref{thm:bisymmetric-two-length-test}の経路長3と21で正則性が決まる。
窓の同型では対応先の四配置にも同じ点数を仮定する。
\end{proposition}
\begin{proof}
補題を各位置の射影変換に使う。テンソル直線の一致は閉条件で、
両対称なら二交換を同時に行う写像が二位置移動の同型である。
その各位置の同定を可逆に延ばせば、一交換ずつのテンソルの等式も特殊化する。
成分の閉点には一般表示の稠密開集合から付値曲線を取ればよい。
一般点の射影同型が代数拡大を要する場合はその曲線を有限被覆してよい。

四元群の各生成元は一般点で二乗1の同変な持上げを持つ。
可逆な射影延長と行列式の付値0によりこの持上げ自身が格子を保つ。
文字冪等元の像は直和因子なので四つの階数1は特殊点でも保たれ、テンソルの文字も保たれる。
両対称表示と同様、有限個のテンソル等式で表す文字付き表示にこの手順を使える。
十六点以上を個々に固定する射影変換は恒等写像である。実際非スカラーな行列の固有空間の和集合には、
上の上界で高々十一点しか載らない。従って同型と周期は四配置間の有限個の点の置換を試し、
対応する射影行列の可逆性と四テンソルの直線の一致を調べて判定できる。
\end{proof}

\begin{proposition}\label{prop:elliptic-component-frames}
各位置の射影点スキームの台が滑らかな非退化楕円四次曲線と有限集合の和である
負成分の窓では、一般の四元群文字枠と、逆捻り後の二つの対称化同定が特殊化する。
そのような点での同型は各鉛筆の四つの特異二次曲面の頂点を置換する。
\end{proposition}
\begin{proof}
一般楕円曲線のHilbert点を固有性で延ばす。特殊曲線は台の唯一の一次元成分を含み、
この被約成分だけで同じHilbert多項式 $4n$ を持つので一致する。
従って曲線族は滑らかである。曲線の $4p\sim\mathcal O(1)$ なる十六点は、
相対Picard上の四倍写像から有限エタールに延びる。点・直線・平面内のこの相異なる点の
個数は高々 $1,2,4$（直線上の三点なら、それを含む全平面が三点の残りの
次数1の束を動かすことになる）である。従って補題と同じ核と像の議論で
二配置間の一般射影同型は延びる。四元群については前命題の文字射影を使う。

捻る前の二交換の合成で四文字の前後二組を交換すると、コサイクル符号の比は
$\mu(i+j,t)\mu(t,i+j)$ である。ただし $t$ は四文字の総次数で、比は二文字 $i,j$ の各文字倍率の積となる。
文字の平行移動による補正も一文字の倍率だけである。よって捻り前の各反対同定は、
格子に可逆な文字倍率を除いて捻り後の位置移動の各射影変換に等しく、可逆に延びる。
逆捻りをして二つの交換式を特殊化すればよい。最後に滑らかな四次完全交叉の二次式は鉛筆をなし、
四つの相異なる特異二次曲面を持つ。その頂点は鉛筆を同時対角化する四頂点そのものである。
\end{proof}

\begin{theorem}\label{thm:klein-bisymmetric-irreducible}
$K=(\mathbb Z/2)^2=\{0,1,2,3\}$ とし、偶二位置・奇二位置をそれぞれ同一視した
正則文字基底を固定する。各々の交換に対称、総文字1、第一切断の色順の階数が
$(2,2,1,1)$ のテンソル全体は既約で次元11である。
従ってこの型の任意の文字付き両対称窓とその非自明コサイクル捻りは、それぞれ上の正・負成分に属する。
負成分のうち前命題の楕円曲線が残る部分では、この逆捻り表示の存在が所属の必要十分条件である。
この部分では零係数の有無によらず同じ経路長3と21の有限正則性判定が使える。
\end{theorem}
\begin{proof}
$M_h(i,k)=w(i,i+h,k,k+h+1)$ とする。二重交換から $M_{h+1}=M_h^{\mathsf t}$。
$M_2=ab^{\mathsf t}\ne0$ は階数1であり、$a,b$ は共通逆倍率を除いて決まる。
一交換ずつの対称性を使うと、$M_0=\left(\begin{smallmatrix}X&U\\V&Y\end{smallmatrix}\right)$ の
$X,Y$ は自由な二行二列、交差ブロックは
\[
 U=\begin{pmatrix}a_2b_0&a_0b_3\\a_1b_2&a_3b_1\end{pmatrix},\qquad
 V=\begin{pmatrix}a_0b_2&a_2b_1\\a_3b_0&a_1b_3\end{pmatrix}.
\]
この表示に残る条件は $\operatorname{rank}M_0=2$ だけである。係数の全軌道と
前命題の符号は \path{computations/verify_klein_bisymmetric_closure.py} で検算する。
七次元の階数1行列の上の十五次元空間で階数2条件の余次元は高々4なので、
解の各成分は少なくとも11次元である（階数多様体のグラフとの交叉次元定理）。
$\Delta=\det U=-\det V\ne0$ なら $X$ は可逆で $Y=VX^{-1}U$、従ってこれは既約な十一変数の開部分である。

$\Delta=0$ のファイバーに限った成分がないことを示す。
$M_0=\binom{L}{R}(C\ D')$ を全列・全行階数2の因子に分けると
$LD'=U$, $RC=V$、共通の $\mathrm{GL}_2$ による四自由度を除く。
$U\ne0$ の二行正方分解 $(L,D')$ は高々4次元である。
一方が可逆なら4次元、双方が階数1なら $U=uv^t$ として
$L=ul^t$, $D'=dv^t$, $l^t d=1$ より3次元である。
$U=0$ なら高々5次元であり、非零な二階数1因子の中間直線と各々の二係数で数えられる。
$V$ も同じで、ファイバーの次元は $4+\mathbf1_{U=0}+\mathbf1_{V=0}$ 以下である。
$U,V\ne0$ の判別式上は基底が高々6次元だから合計高々10である。
一方のブロック全体が零なら四成分の各々は相異なる $a_i$ と相異なる $b_j$ の積である。
従って少なくとも四座標が零で、基底は逆倍率を除き高々3次元、合計高々9となる。
これで既約性は疎な全ての台にも及ぶ。

一般の固定枠ではこの既約開部分が定理\ref{thm:rectangle-unrestricted-components}の
正の四変数表示に達する。同変同定を動かす六自由度と固定基底軌道の十三自由度を
考慮すれば $11+6-13=4$ である。捻りは固定色のテンソル空間の可逆線型変換だから
負成分への包含も従う。逆向きに前命題で延びた枠と二同定を取ると、
各色の階数は一般値から増えず、和が6のままなので個々の階数は落ちない。
逆捻りの各隣接切断は色別の左右対角倍率を除いて元の切断に等しい
（交差する二対の符号は色ごとの定数である）。
よってこの型の表示が復元される。捻りは経路の次元と左右Koszul・AS性を保つので
両対称の有限判定が全ての表示に適用できる。
\end{proof}

\begin{corollary}\label{cor:negative-component-period}
負成分の幾何学的一般点の最小周期は2、指定始点の射影安定群は四元群である。
滑らかな楕円四次が各位置の唯一の曲線として残る軌跡では、任意の二点の同型、
各点の周期と安定群を有限個の置換と整数格子で判定できる。
\end{corollary}
\begin{proof}
各鉛筆の四頂点の枠でテンソルを書く。始点の移動を指定し、四位置の置換で台を合わせた後、
対応する非零係数の比を $r_q$ とする。$E$ を台の各語 $q$ を行とし
$E_{q,(p,i)}=[q_p=i]$ とした整数行列とする。対角倍率の存在は
$\prod_q r_q^{v_q}=1$ が $\ker_{\mathbb Z} E^{\mathsf t}$ の基底の全てで成り立つことと
同値である。これは代数閉体上のトーラス準同型の像の文字による表示で、零係数のある台にもそのまま使える。
解は一つあれば同じ対角安定群の剰余類であり、各位置のスカラーを除けば窓の射影安定群を得る。

一般表示では対角鉛筆の四頂点は元の文字線である。非零台は
$i+j+k+l=0$ の64語全体で、これを保つ置換は $i_p\mapsto\phi(i_p)+u_p$、
$\phi\in\mathrm{GL}(2,\mathbb F_2)$, $\sum_p u_p=0$ に限る。実際、各置換の零での値を引き、
二位置に等しい文字を置けば四写像は一致し、三位置に $i,j,i+j$ を置けば共通写像は加法的である。
$E$ のSmith因子は $1^{11},2,2$ である。上の有理標本の剰余1009で各移動につき384通りを調べ、
移動 $0,1,2,3$ の順に残るのは $1,0,1,0$ 通りである。
全ての係数が単元であるこの点で失敗する単項式の等式は、関数体上の恒等式ではない。
恒等置換と既に示した二位置移動は一般点でも成立し、対角安定群の有限部分はSmith因子から位数4である。
\path{computations/verify_negative_component_period.py} が飽和整数核を使って検算する。
\end{proof}

定理\ref{thm:rectangle-unrestricted-components}の三つの水平標本と本節の $A_1=A_0$ の
正標本では、四位置の全てで十六点条件を適用できる。
長さ20の有限自由点環のアフィン四座標を $a_0,\ldots,a_3$ とすると、
$\sum 2^i a_i$ の乗法の特性多項式の平方自由部分の次数は、剰余1009で標本順に $18,19,18,20$ である。
これだけの相異なる値は標数零にも持ち上がる。
\path{computations/verify_rectangle_stable_configurations.py} が全十六配置で検算する。
